\hnheader[Der 0/1-Verlust ist nicht optimierbar -- also nutzen wir konvexe Ersatzverluste.][Margin $z=y\,\theta\T x$: $z>0$ heißt richtig klassifiziert!]{Verlust-}{funktionen}{Konvexe Ersatzverluste (Surrogate Losses)}
\hnsrc{extra-notes/loss-functions.pdf und extra-notes/representer-function.pdf (J.\ Duchi)}

\begin{hngrid}
%
\begin{hnp}[raster multicolumn=2,acc=hnorange]{1}{Der Margin}
Für $y\in\{-1,+1\}$ und $\hat y=\theta\T x$ ist $z=y\,\theta\T x$ der \hlt{Margin}: $z>0$ richtig, $z<0$ falsch; $|z|$ misst die Sicherheit. Empirisches Risiko: $J(\theta)=\frac1n\sum_i\phi(y^{(i)}\theta\T x^{(i)})$.
\end{hnp}
%
\begin{hnp}[raster multicolumn=2,acc=hnpurple]{2}{Problem: 0/1-Verlust}
$\phi_{01}(z)=\I\{z\le0\}$ zählt Fehler, ist aber unstetig, nicht konvex und NP-hart zu minimieren.\\
Wunsch: $\phi$ konvex, stetig, $\phi(z)\to0$ für $z\to\infty$, $\phi(z)\to\infty$ für $z\to-\infty$.
\end{hnp}
%
\begin{hnp}[raster multicolumn=2,acc=hnteal]{3}{Skizze: Verluste über $z$}
\centering
\begin{tikzpicture}[x=.95cm,y=.66cm,>=Stealth]
  \draw[->,hnnavy] (-2.2,0)--(3.2,0) node[right,font=\tiny]{$z$};
  \draw[->,hnnavy] (0,-.1)--(0,4.3);
  \draw[hnnavy,thick] (-2.2,1)--(0,1); \draw[hnnavy,thick] (0,0)--(3,0); \draw[hnnavy,dotted] (0,0)--(0,1);
  \draw[hnorange,thick,domain=-2.2:3,samples=50] plot(\x,{ln(1+exp(-\x))/ln(2)});
  \draw[hnpurple,thick] (-2.2,3.2)--(1,0)--(3,0);
  \draw[hnrose,thick,domain=-1.1:3,samples=50] plot(\x,{exp(-\x)});
  \node[font=\tiny,hnnavy] at (-1.6,1.4) {0/1};
  \node[font=\tiny,hnorange] at (-1.7,2.55) {logistisch};
  \node[font=\tiny,hnpurple] at (-1.3,3.5) {Hinge};
  \node[font=\tiny,hnrose] at (-.35,3.7) {exp};
\end{tikzpicture}
\end{hnp}
%
\begin{hnp}[raster multicolumn=3,acc=hnnavy]{4}{Drei Verluste, drei Verfahren}
\renewcommand{\arraystretch}{1.3}
\begin{tabular}{@{}p{2.0cm}p{2.8cm}p{3.6cm}@{}}
\hline
\textbf{Verlust}&$\phi(z)$&\textbf{Verfahren}\\\hline
Logistisch&$\log(1+e^{-z})$&Logistische Regression\\
Hinge&$\max(1-z,0)$&Support Vector Machine\\
Exponentiell&$e^{-z}$&AdaBoost\\\hline
\end{tabular}
\\[3pt]
\hlt{Verbindendes Prinzip:} alle minimieren $\frac1n\sum\phi(y\theta\T x)$ -- nur $\phi$ unterscheidet sich; alle sind konvexe obere Schranken des 0/1-Verlusts.
\end{hnp}
%
\begin{hnp}[raster multicolumn=3,acc=hngreen]{5}{Generisches Training}
\begin{pcode}[ERM mit Ersatzverlust]
\kw{init} $\theta\leftarrow0$\\
\kw{repeat}: wähle Beispiel $(x_i,y_i)$\\
\hii $z\leftarrow y_i\,\theta^{\top}x_i$\\
\hii $g\leftarrow\phi'(z)\cdot y_i\,x_i$\\
\hii $\theta\leftarrow\theta-\alpha\,g$\\
\cm{\# $\alpha$: Lernrate, $\phi'$: Ableitung; logistisch $-1/(1+e^{z})$}\\
\cm{\# Hinge: $-1$ falls $z<1$, sonst $0$ (Subgradient)}\\
\cm{\# exp: $-e^{-z}$}
\end{pcode}
\end{hnp}
%
\begin{hnex}[raster multicolumn=3]{Beispiel: Verluste bei drei Margins}
\renewcommand{\arraystretch}{1.25}
\begin{tabular}{@{}c|cccc@{}}
$z$&0/1&logist.&Hinge&exp\\\hline
$3$&$0$&$0.05$&$0$&$0.05$\\
$0.5$&$0$&$0.47$&$0.5$&$0.61$\\
$-1$&$1$&$1.31$&$2$&$2.72$\\
\end{tabular}\\[2pt]
Selbst \emph{richtig} klassifizierte Punkte mit kleinem Margin ($z=0.5$) werden noch bestraft $\Rightarrow$ Ansporn, den Margin zu vergrößern.
\end{hnex}
%
\begin{hnp}[raster multicolumn=3,acc=hnrose]{6}{Verallgemeinerte Verluste}
Allgemein: $L(\theta\T x,y)$ ($z$: Score, $k$: Anzahl Klassen). Regression: $\frac12(z-y)^2$. Binär: $\log(1+e^{-yz})$. Multiklasse mit $\Theta=[\theta_1\cdots\theta_k]$:
\[ L(\Theta\T x,y)=\log\sum_{i=1}^{k}e^{x\T(\theta_i-\theta_y)} \]
(Cross-Entropy/Softmax; $\theta_y$: Parameter der wahren Klasse, $\Theta$: alle Klassenparameter als Spalten). Mit konvexem $L$ und $\ell_2$-Regularisierung gilt das Representer-Theorem.
\end{hnp}
%
\begin{hnp}[raster multicolumn=2,acc=hngreen]{7}{Wann welcher Verlust?}
\begin{itemize}
\item Wahrscheinlichkeiten nötig $\to$ logistisch
\item Margin/Sparsität $\to$ Hinge
\item Boosting $\to$ exponentiell (ausreißeranfällig)
\end{itemize}
\end{hnp}
%
\begin{hnp}[raster multicolumn=2,acc=hnorange]{8}{Key Takeaways}
\begin{hnchk}
\item Margin $z=y\theta\T x$
\item 0/1 nicht optimierbar $\Rightarrow$ Surrogat
\item logistisch = LogReg, Hinge = SVM, exp = AdaBoost
\end{hnchk}
\end{hnp}
%
\begin{hnp}[raster multicolumn=2,acc=hnteal]{9}{Querverweise}
\begin{itemize}
\item Logistische Regression, SVM, Ensembles
\item Praktische Aufgabe: \texttt{p01\_lr.py} (Aufgabenblatt 2) nutzt genau diese $\pm1$-Notation
\end{itemize}
\end{hnp}
%
\begin{hnp}[raster multicolumn=6,acc=hnpurple,colback=hnlilac!30]{?}{Selbsttest: kannst du das erklären?}
\hnquiz{Warum nicht den 0/1-Verlust minimieren?}{Unstetig, nichtkonvex, NP-hart.}{Welcher Verlust gehört zu SVM, AdaBoost, LogReg?}{Hinge, exponentiell, logistisch.}{Was ist der Margin?}{$z=y\theta\T x$: $z>0$ richtig, $|z|$ Sicherheit.}
\end{hnp}
%
\begin{hnbanner}[raster multicolumn=6]
„Optimiere, was du optimieren kannst -- und hoffe, dass es mit dem zusammenfällt, was du willst.“
\end{hnbanner}
\end{hngrid}
